\section{Conclusion}
\label{sec:conclusion}

This paper shows that long-horizon game agents benefit from coupling structured
control with structured memory. \textbf{\method{}} uses adaptive dual control to reserve
expensive strategic reasoning for planning and recovery, while leaving routine
execution to a lightweight reactive controller. Together with Hierarchical
Memory, this design improves success, cost, and recovery on Lite-100 and also
transfers to RDR2 under different visual-domain and action-interface assumptions.
Ablations further show that \textbf{event triggering}, \textbf{reactive
override}, and \textbf{controller-specific memory} are complementary rather than
interchangeable parts of the framework.

\paragraph{Limitations and future work.}
\label{sec:limitations}
\method{} can still misroute control when progress is visually subtle, failure
signals are noisy, when a game state changes without large screen differences,
or when retrieved memory becomes stale. Its evidence is strongest on Lite-100
and selected RDR2 tasks, so broader game styles, lower-level interfaces, and
longer deployments remain important future tests. The current system uses visual
observations and action history, while audio cues and richer multimodal feedback
are not modeled. Although \method{} reduces unnecessary large-model calls,
recovery and replanning still depend on explicit model invocations, leaving cost
and latency as practical constraints for deployment. Future work should learn
event-triggering and memory-writing policies from interaction data, add
uncertainty-aware retrieval, incorporate additional task signals, and study
whether training or demonstrations can improve the lightweight controller
without sacrificing adaptability.
